\documentclass[11pt,a4paper]{article}
% Replace these two fields before submission; no identity has been inferred.
\newcommand{\StudentName}{[YOUR NAME]}
\newcommand{\StudentID}{[YOUR STUDENT ID]}

\usepackage{assignment}

\begin{document}
\makeassignmenttitle

\question{1}
The two lines have the common direction $\mathbf d=(4,2,5)^\top$.
For a point $\mathbf a+\lambda\mathbf d$ on either line,
\[
 \lim_{\lambda\to\infty}
 \left(\frac{a_X+4\lambda}{a_Z+5\lambda},
       \frac{a_Y+2\lambda}{a_Z+5\lambda}\right)
 =\left(\frac45,\frac25\right).
\]
Thus both lines have the same vanishing point
$\answer{\mathbf v=(4/5,\,2/5)^\top=(0.8,\,0.4)^\top}$.
Its homogeneous image coordinates are $(4,2,5)^\top$; the different line origins
do not affect this limit.

\question{2}
The plane has normal $\mathbf n=(6,3,2)^\top$.
Every direction $\mathbf d$ parallel to it satisfies $\mathbf n^\top\mathbf d=0$.
Under normalized perspective projection, its vanishing point is
$\mathbf v\sim(d_X,d_Y,d_Z)^\top$. Therefore all these points satisfy
\[
 \mathbf n^\top\mathbf v=0
 \quad\Longrightarrow\quad
 \answer{6x+3y+2=0},\qquad
 \answer{\boldsymbol\ell_\infty\sim(6,3,2)^\top}.
\]
The offset $17$ determines the plane's position, whereas its vanishing line
depends only on its orientation. Directions with $d_Z=0$ are included as
homogeneous points at infinity.

\question{3}
For $12X+8Y+28Z=23$, a normal is
\[
 \mathbf n=(12,8,28)^\top=4(3,2,7)^\top.
\]
If $\mathbf u$ and $\mathbf v$ lie on the plane, subtracting their plane
equations gives $\mathbf n^\top(\mathbf u-\mathbf v)=0$.
Consequently $(3,2,7)^\top$ is orthogonal to every direction in the plane,
so it is perpendicular to the plane.

\question{4}
The shared translation $\mathbf t=(-3,2,1)^\top$ cancels:
$\mathbf x_2-\mathbf y_2=R(\mathbf x_1-\mathbf y_1)$.
Writing $\mathbf d=\mathbf x_1-\mathbf y_1$ and using $R^\top R=I$,
\[
 \|\mathbf x_2-\mathbf y_2\|_2^2
 =\mathbf d^\top R^\top R\mathbf d
 =\mathbf d^\top\mathbf d
 =\|\mathbf x_1-\mathbf y_1\|_2^2.
\]
Both norms are nonnegative; taking square roots proves
$\answer{\|\mathbf x_2-\mathbf y_2\|_2=\|\mathbf x_1-\mathbf y_1\|_2}$.

\clearpage
\question{5}
Let $\mathbf X_i=(X_i,Y_i,Z_i,1)^\top$ and denote the rows of $P$ by
$\mathbf p_1^\top,\mathbf p_2^\top,\mathbf p_3^\top$.
Each correspondence imposes the two linear constraints
\[
 (\mathbf p_1^\top-x_i\mathbf p_3^\top)\mathbf X_i=0,
 \qquad
 (\mathbf p_2^\top-y_i\mathbf p_3^\top)\mathbf X_i=0.
\]
Stacking the six pairs produces $A\mathbf p=0$, where $A$ is $12\times12$
and $\mathbf p$ contains the rows of $P$ in order. Its rank is $11$, so the
solution is unique up to a nonzero scale. The exact solution has representative
\[
 P_0=\begin{bmatrix}
 180&0&0&280\\0&180&0&220\\0&0&0&1
 \end{bmatrix},\qquad
 \|P_0\|_F=\sqrt{180^2+280^2+180^2+220^2+1}=\sqrt{191601}.
\]
The requested normalization is
\[
 \answer{P=\frac{1}{\sqrt{191601}}
 \begin{bmatrix}180&0&0&280\\0&180&0&220\\0&0&0&1\end{bmatrix}}.
\]
This camera gives $x=180X+280$, $y=180Y+220$, independent of $Z$:
it is an affine camera. In particular, the fifth point gives $(370,310)$
and the sixth gives $(370,220)$. A finite pinhole intrinsic matrix is not
an additional assumption of this computation. The matrices $P$ and $-P$
have the same Frobenius norm and project to the same image points; the
implementation fixes the sign by making the largest-magnitude entry positive.

\textbf{Implementation.}
\codepath{experients/assignment1/q5_camera.py}, reproduced in Appendix A.
The code centers and scales the 3D and 2D points, takes the last right singular
vector of the normalized design matrix, and denormalizes with
$P=T_{2D}^{-1}P_nT_{3D}$ before applying the Frobenius normalization.
The following is captured from the script's actual \texttt{print} output.
\lstinputlisting[style=output]{../experients/assignment1/results/q5_camera.txt}

\clearpage
\question{6}
Let $\mathbf X_h=(X,Y,Z,W)^\top$. For each camera, stack the rows
$x_jP^{(j)}_3-P^{(j)}_1$ and $y_jP^{(j)}_3-P^{(j)}_2$.
The supplied matrices and pixels give
\[
 A=\begin{bmatrix}
 -900&0&0&0\\0&-900&0&0\\-900&0&-15&6\\0&-900&0&0
 \end{bmatrix},\qquad A\mathbf X_h=0.
\]
The first two equations imply $X=Y=0$; the third then gives
$-15Z+6W=0$. Choosing $W=1$ yields
\[
 \answer{\mathbf X=(0,\,0,\,0.4)^\top},\qquad
 \mathbf X_h=(0,\,0,\,0.4,\,1)^\top.
\]
The fourth column of $P^{(2)}$ is already expressed in projection-matrix
units. Indeed, $P^{(2)}=K[I\mid\mathbf t]$ has
$\mathbf t=(-6/900,0,0)^\top$ and camera center
$\mathbf C_2=(6/900,0,0)^\top$.
Hence the physical baseline is $B=6/900=1/150$ in the given world units,
and the disparity check is
\[
 Z=\frac{fB}{x_1-x_2}=\frac{900(6/900)}{400-385}=0.4.
\]
Both camera depths are positive. Direct reprojection confirms
\[
 \pi(P^{(1)}\mathbf X_h)=(400,250),\qquad
 \pi(P^{(2)}\mathbf X_h)=(385,250),
 \quad \pi(a,b,c)=(a/c,b/c).
\]

\textbf{Implementation.}
\codepath{experients/assignment1/q6_triangulation.py}, reproduced in Appendix B,
solves the homogeneous system by SVD and divides by the last coordinate.
Its captured \texttt{print} output is:
\lstinputlisting[style=output]{../experients/assignment1/results/q6_triangulation.txt}

\clearpage
\question{7}
With the convention $\mathbf X_2=R\mathbf X_1+\mathbf t$, corresponding
calibrated image points obey $\mathbf p_2^\top E\mathbf p_1=0$, with
\[
 E=[\mathbf t]_\times R,\qquad
 [\mathbf t]_\times=\begin{bmatrix}
 0&-t_z&t_y\\t_z&0&-t_x\\-t_y&t_x&0
 \end{bmatrix}.
\]
For $R=I$ and $\mathbf t=(3,0,0)^\top$,
\[
 \answer{E=\begin{bmatrix}0&0&0\\0&0&-3\\0&3&0\end{bmatrix}}.
\]
It is defined up to nonzero scale. Since $K_1=K_2=I$, the fundamental
matrix $F=K_2^{-\top}EK_1^{-1}$ is also $E$.

\question{8}
The epipolar line in image 2 is $\boldsymbol\ell_2=E\mathbf p_1$:
\[
 \boldsymbol\ell_2=
 \begin{bmatrix}0&0&0\\0&0&-3\\0&3&0\end{bmatrix}
 \begin{bmatrix}280\\200\\1\end{bmatrix}
 =\begin{bmatrix}0\\-3\\600\end{bmatrix}.
\]
Thus $\answer{\boldsymbol\ell_2\sim(0,1,-200)^\top}$, or equivalently
$\answer{y_2=200}$, since $0x_2-3y_2+600=0$.

\question{9}
First multiply and then divide by the homogeneous image coordinate:
\[
 P\mathbf X=\begin{bmatrix}3\\1\\2\end{bmatrix}
 \quad\Longrightarrow\quad
 \widehat{\mathbf x}=(1.5,\,0.5)^\top.
\]
Using the supplied projection matrix and measurement directly, in the pixel
units requested by the question, the Euclidean reprojection error is
\begin{align*}
 e&=\|\widehat{\mathbf x}-\mathbf x_m\|_2
   =\sqrt{(1.5-280)^2+(0.5-200)^2}\\
  &=\sqrt{278.5^2+199.5^2}
   =\sqrt{117362.5}
   \approx\answer{342.5821\ \text{px}}.
\end{align*}
No extra intrinsic matrix or coordinate conversion is specified.

\question{10}
Unit-scale orthographic projection along $+Z$ discards the depth coordinate:
\[
 \mathbf x=\begin{bmatrix}1&0&0\\0&1&0\end{bmatrix}
 \begin{bmatrix}5\\4\\8\end{bmatrix}
 =\answer{\begin{bmatrix}5\\4\end{bmatrix}}.
\]

\clearpage
\section{Appendix A: complete Python code for Question 5}
File: \codepath{experients/assignment1/q5_camera.py}.
This is the same source file that generated the output in Question 5.
\lstinputlisting[style=python]{../experients/assignment1/q5_camera.py}

\section{Appendix B: complete Python code for Question 6}
File: \codepath{experients/assignment1/q6_triangulation.py}.
This is the same source file that generated the output in Question 6.
\lstinputlisting[style=python]{../experients/assignment1/q6_triangulation.py}

\clearpage
\section{Appendix C: verification and reproducibility}
The experiments require Python, NumPy and SymPy. From the repository root:
\begin{lstlisting}[style=output]
python -m pip install -r experients/assignment1/requirements.txt
python experients/assignment1/run.py
powershell -NoProfile -File report/build.ps1
\end{lstlisting}
The build script reruns the experiments and compiles this PDF with XeLaTeX.
The source for the independent checks is
\codepath{experients/assignment1/verify.py}.
It uses exact rational arithmetic for the camera's rank and null vector,
exact reprojection for triangulation, and independent projected points to
check the epipolar convention. Random checks use seed 1050; the numerical
checks supplement the proofs above. Its actual printed results are:
\lstinputlisting[style=output]{../experients/assignment1/results/verify.txt}

\subsection{Recorded execution environment}
\begin{tabular}{ll}
\toprule
Component & Recorded value \\
\midrule
Python & 3.10.16 \\
NumPy & 2.2.6 \\
SymPy & 1.13.3 \\
Manifest SHA-256 (first 16 hex) & \texttt{8d96f70449b47602} \\
\bottomrule
\end{tabular}


The complete SHA-256 hashes of the assignment inputs, code, report sources
and printed outputs are in
\codepath{experients/assignment1/results/manifest.json}.
The file \codepath{experients/assignment1/inputs.md} transcribes the numerical
inputs and records the submission requirements. Original task sources are
\texttt{assignment1\_CS1050.pdf} and the accompanying submission notice.
Rounded matrix displays may contain signed zeros; validation uses the
unrounded arrays. The final PDF includes both complete question-specific
Python files and their printed results, so it can be submitted as one file.

\end{document}
